% GENERATED FILE. Edit canonical lessons, then run npm run generate:books.
% canonical-source-hash: fnv1a64:9e1659026bafb01b
% canonical-lessons: AR-C482-nasim, AR-C482-isar, AR-C482-asafa, AR-C482-ishtadda, AR-C482-khatar

\chapter{Wind and storms}
\label{ch:ar-wind-and-storms}
\hlchaptermodality{482}

\begin{chapteropening}
\textbf{By the end of this chapter:} \emph{I can say a breeze, a hurricane, to blow hard, to grow stronger and danger.}

Complete the last lesson of chapter 482: I can say a breeze, a hurricane, to blow hard, to grow stronger and danger.
\end{chapteropening}

\section[\texorpdfstring{\ar{نسيم} (nasīm)}{نسيم (nasīm)}]{\ar{نسيم} (nasīm) — a breeze}

\label{lesson:AR-C482-nasim}



\begin{quote}
Before the new one: say the Arabic for to pour down, then the Arabic for to clear up.
\end{quote}

\subsection*{You'll want to know: \ar{نسيم}}
\textbf{\ar{نسيم}} — \emph{nasīm} — \textquotedblleft{}a breeze\textquotedblright{}.

\begin{quote}
\ar{شعرنا} \ar{بنسيم} \ar{لطيف} \ar{على} \ar{الشرفة}.
\end{quote}

\begin{quote}
\emph{shaʿarnā bi-nasīm laṭīf ʿalā ash-shurfa.}
\end{quote}

\begin{quote}
We felt a pleasant breeze on the balcony.
\end{quote}

\begin{cousinweb}
\textbf{\ar{نسيم}} (\emph{nasīm}) is built on the native Arabic root \emph{n-s-m}, \textquotedblleft{}to blow gently; breath\textquotedblright{} (Lane).
\end{cousinweb}

\begin{grammarlens}[title={What you've built}]
One more word for getting around.
\end{grammarlens}

\subsection*{Your turn}
\begin{itemize}
\raggedright
  \item \emph{Say it:} \emph{nasīm}
  \item \emph{Say it:} \emph{nasīm}, once more
  \item \emph{Say it:} \emph{shaʿarnā bi-nasīm laṭīf ʿalā ash-shurfa.} — We felt a pleasant breeze on the balcony.
  \item \emph{Recall:} say the Arabic for to pour down, then the Arabic for to clear up, then say \emph{nasīm} again
\end{itemize}

\begin{tcolorbox}[breakable,colback=teal!4,colframe=teal!35!black,title={Before you move on}]
What does \ar{نسيم} mean? (\textquotedblleft{}A breeze\textquotedblright{}.) Say it once more.
\end{tcolorbox}
\section[\texorpdfstring{\ar{إعصار} (iʿṣār)}{إعصار (iʿṣār)}]{\ar{إعصار} (iʿṣār) — a hurricane; a whirlwind}

\label{lesson:AR-C482-isar}



\begin{quote}
Before the new one: say the Arabic for to clear up, then the Arabic for a breeze.
\end{quote}

\subsection*{You'll want to know: \ar{إعصار}}
\textbf{\ar{إعصار}} — \emph{iʿṣār} — \textquotedblleft{}a hurricane; a whirlwind\textquotedblright{}.

\begin{quote}
\ar{وصل} \ar{إعصار} \ar{قوي} \ar{إلى} \ar{الساحل}.
\end{quote}

\begin{quote}
\emph{waṣala iʿṣār qawī ilā as-sāḥil.}
\end{quote}

\begin{quote}
A strong hurricane reached the coast.
\end{quote}

\begin{cousinweb}
\textbf{\ar{إعصار}} (\emph{iʿṣār}) is built on the native Arabic root \emph{ʿ-ṣ-r}, \textquotedblleft{}to press, to wring\textquotedblright{} (Lane). Lane gives \emph{iʿṣār} as a whirlwind that raises the dust.
\end{cousinweb}

\begin{grammarlens}[title={What you've built}]
One more word for getting around.
\end{grammarlens}

\subsection*{Your turn}
\begin{itemize}
\raggedright
  \item \emph{Say it:} \emph{iʿṣār}
  \item \emph{Say it:} \emph{iʿṣār}, once more
  \item \emph{Say it:} \emph{waṣala iʿṣār qawī ilā as-sāḥil.} — A strong hurricane reached the coast.
  \item \emph{Recall:} say the Arabic for to clear up, then the Arabic for a breeze, then say \emph{iʿṣār} again
\end{itemize}

\begin{tcolorbox}[breakable,colback=teal!4,colframe=teal!35!black,title={Before you move on}]
What does \ar{إعصار} mean? (\textquotedblleft{}A hurricane; a whirlwind\textquotedblright{}.) Say it once more.
\end{tcolorbox}
\section[\texorpdfstring{\ar{عصف} (ʿaṣafa)}{عصف (ʿaṣafa)}]{\ar{عصف} (ʿaṣafa) — to blow hard (of wind)}

\label{lesson:AR-C482-asafa}



\begin{quote}
Before the new one: say the Arabic for a breeze, then the Arabic for a hurricane.
\end{quote}

\subsection*{You'll want to know: \ar{عصف}}
\textbf{\ar{عصف}} — \emph{ʿaṣafa} — \textquotedblleft{}to blow hard (of wind)\textquotedblright{}.

\begin{quote}
\ar{عصفت} \ar{الريح} \ar{في} \ar{الليل}.
\end{quote}

\begin{quote}
\emph{ʿaṣafat ar-rīḥ fī al-layl.}
\end{quote}

\begin{quote}
The wind blew hard in the night.
\end{quote}

Present: \ar{يعصف}, \emph{yaʿṣif}.

\begin{cousinweb}
\textbf{\ar{عصف}} (\emph{ʿaṣafa}) is built on the native Arabic root \emph{ʿ-ṣ-f}, \textquotedblleft{}to blow violently (of wind)\textquotedblright{} (Lane). It is the plain verb; \emph{ʿāṣifa}, \textquotedblleft{}a storm\textquotedblright{}, is its participle.
\end{cousinweb}

\begin{grammarlens}[title={What you've built}]
One more word for getting around.
\end{grammarlens}

\subsection*{Your turn}
\begin{itemize}
\raggedright
  \item \emph{Say it:} \emph{ʿaṣafa}
  \item \emph{Say it:} \emph{ʿaṣafa}, once more
  \item \emph{Say it:} \emph{ʿaṣafat ar-rīḥ fī al-layl.} — The wind blew hard in the night.
  \item \emph{Recall:} say the Arabic for a breeze, then the Arabic for a hurricane, then say \emph{ʿaṣafa} again
\end{itemize}

\begin{tcolorbox}[breakable,colback=teal!4,colframe=teal!35!black,title={Before you move on}]
What does \ar{عصف} mean? (\textquotedblleft{}To blow hard (of wind)\textquotedblright{}.) Say it once more.
\end{tcolorbox}
\section[\texorpdfstring{\ar{اشتد} (ishtadda)}{اشتد (ishtadda)}]{\ar{اشتد} (ishtadda) — to grow stronger; to intensify}

\label{lesson:AR-C482-ishtadda}



\begin{quote}
Before the new one: say the Arabic for a hurricane, then the Arabic for to blow hard.
\end{quote}

\subsection*{You'll want to know: \ar{اشتد}}
\textbf{\ar{اشتد}} — \emph{ishtadda} — \textquotedblleft{}to grow stronger; to intensify\textquotedblright{}.

\begin{quote}
\ar{اشتدت} \ar{العاصفة} \ar{في} \ar{المساء}.
\end{quote}

\begin{quote}
\emph{ishtaddat al-ʿāṣifa fī al-masāʾ.}
\end{quote}

\begin{quote}
The storm grew stronger in the evening.
\end{quote}

Present: \ar{يشتد}, \emph{yashtadd}.

\begin{cousinweb}
\textbf{\ar{اشتد}} (\emph{ishtadda}) is built on the native Arabic root \emph{sh-d-d}, \textquotedblleft{}to be firm; to bind tightly\textquotedblright{} (Lane). It is Form VIII.
\end{cousinweb}

\begin{grammarlens}[title={What you've built}]
One more word for getting around.
\end{grammarlens}

\subsection*{Your turn}
\begin{itemize}
\raggedright
  \item \emph{Say it:} \emph{ishtadda}
  \item \emph{Say it:} \emph{ishtadda}, once more
  \item \emph{Say it:} \emph{ishtaddat al-ʿāṣifa fī al-masāʾ.} — The storm grew stronger in the evening.
  \item \emph{Recall:} say the Arabic for a hurricane, then the Arabic for to blow hard, then say \emph{ishtadda} again
\end{itemize}

\begin{tcolorbox}[breakable,colback=teal!4,colframe=teal!35!black,title={Before you move on}]
What does \ar{اشتد} mean? (\textquotedblleft{}To grow stronger; to intensify\textquotedblright{}.) Say it once more.
\end{tcolorbox}
\section[\texorpdfstring{\ar{خطر} (khaṭar)}{خطر (khaṭar)}]{\ar{خطر} (khaṭar) — danger}

\label{lesson:AR-C482-khatar}



\begin{quote}
Before the new one: say the Arabic for to blow hard, then the Arabic for to grow stronger.
\end{quote}

\subsection*{You'll want to know: \ar{خطر}}
\textbf{\ar{خطر}} — \emph{khaṭar} — \textquotedblleft{}danger\textquotedblright{}.

\begin{quote}
\ar{كان} \ar{الخطر} \ar{كبيرا} \ar{على} \ar{الساحل}.
\end{quote}

\begin{quote}
\emph{kāna al-khaṭar kabīran ʿalā as-sāḥil.}
\end{quote}

\begin{quote}
The danger was great on the coast.
\end{quote}

Plural: \ar{أخطار}, \emph{akhṭār}.

\begin{cousinweb}
\textbf{\ar{خطر}} (\emph{khaṭar}) is built on the native Arabic root \emph{kh-ṭ-r}, \textquotedblleft{}peril; to be of great weight\textquotedblright{} (Lane). It shares its root with \emph{khaṭīr}, \textquotedblleft{}dangerous\textquotedblright{}.
\end{cousinweb}

\begin{grammarlens}[title={What you've built}]
One more word for getting around.
\end{grammarlens}

\subsection*{Your turn}
\begin{itemize}
\raggedright
  \item \emph{Say it:} \emph{khaṭar}
  \item \emph{Say it:} \emph{khaṭar}, once more
  \item \emph{Say it:} \emph{kāna al-khaṭar kabīran ʿalā as-sāḥil.} — The danger was great on the coast.
  \item \emph{Recall:} say the Arabic for to blow hard, then the Arabic for to grow stronger, then say \emph{khaṭar} again
\end{itemize}

\begin{tcolorbox}[breakable,colback=teal!4,colframe=teal!35!black,title={Before you move on}]
What does \ar{خطر} mean? (\textquotedblleft{}Danger\textquotedblright{}.) Say it once more.
\end{tcolorbox}
